\subsection*{Cue retention, finite rescue and auxiliary predictive readouts}
Additional controls in the same reduced ecology distinguish remembered content from the procedure used to revise it. A complementary controlled-density model provides repeated resource pulses, a finite common energy reserve, density transport, uptake, growth and turnover. Its actuator moves the whole density distribution towards designated sites; this is a deliberately reduced model, not a locally actuated Lenia organism. A cue-conditioned two-by-two value table learns from actual consumed resource, holds the departure cue and pays for memory access. It supplies a feasibility comparison for remembered prediction, not evidence for the spatial rule's performance in a body.

After thirty-two training cycles, the value table is frozen for sixteen test cycles. In the predictive environment, the history policy captures a mean 0.9955 of available test input across three specified shapes and three seeds, compared with 0.3454 for the no-history bandit. Mean test changes in mass and energy are both positive (Table~\ref{tab:ecology}). These are descriptive means across nine construction--seed cells, not nine independently evolved bodies. A policy given additional resource information approaches complete capture and supplies an opportunity ceiling rather than an information-matched competitor.

\begin{table}[htbp]\centering\small
\caption{Frozen deployment in the separate controlled-density model. Entries are arithmetic means over three specified shapes and three seeds, not independent evolutionary replicates. $\Delta M$ and $\Delta E$ span sixteen test cycles. The oracle receives additional information. The reversal occurs during training. Full results include five further comparison policies.}
\label{tab:ecology}
\begin{tabular}{llrrr}\toprule
World & Policy & Capture fraction & $\Delta M$ & $\Delta E$\\\midrule
Predictive & Cue history & 0.9955 & +9.356 & +9.840\\
Predictive & No-history bandit & 0.3454 & -3.433 & -6.050\\
Predictive & Oracle & 0.9988 & +9.075 & +5.070\\
\addlinespace
Reversal & Cue history & 0.4726 & +3.003 & -17.275\\
Reversal & No-history bandit & 0.1782 & -7.695 & -5.170\\
Reversal & Oracle & 0.9988 & +9.075 & +5.073\\
\addlinespace
Independent & Cue history & 0.2564 & -7.649 & -11.390\\
Independent & No-history bandit & 0.5270 & +3.944 & -9.419\\
Independent & Oracle & 0.9987 & +9.067 & +4.969\\
\bottomrule\end{tabular}
\end{table}


The result changes with the environment. Reversal at cycle 24 occurs during training, before the sixteen frozen deployment cycles begin at cycle 32. Mean deployment capture falls to 0.4726 and energy changes by $-17.275$, despite positive mass change. With independent cues, capture is 0.2564 versus 0.5920 for fixed right-site residence, and five of nine history-policy cases end at exactly zero mass. These descriptive means retain all three seed tapes and their prescribed shapes; the fixed policy's 59\% is not a population success probability. A phase clock, aggregate resource observations and a prescribed return route remain available to the controller.

Matched-state continuations identify what the remembered cue contributes (Table~\ref{tab:ecology-continuations}). From each of nine learned stable-task endpoints, two six-step docking histories with opposite cues leave density, resource, energy and the current observation equal; all nine diagnostic pairs choose opposite release actions. Keeping the learned value table fixed, resetting the cue to zero at every release reduces subsequent capture from 1.000 to 0.542 and final energy from 50.828 to 22.382. Resetting selects the non-positive-cue row of that table, so this is a directed content intervention, not erasure of the learned mapping or unbiased guessing. Intact memory incurs more realized expenditure and still loses 4.386 energy units over the continuation. Useful cue content therefore improves a bodily outcome without guaranteeing sustained energetic gain.

Reopening learning after failed deployment provides a separate, finite recovery test. Each actual low-energy reversal endpoint is split into a continuously frozen future and a future with sixteen cycles of updating and exploration followed by sixteen frozen evaluation cycles. The already reversed mapping is retained; no new reversal is introduced at this branch. Six cases sharing two seeds recover nearly complete capture, while all three cases of the remaining seed die; all nine frozen controls die. Averaged over all nine cases, recovery-branch capture is 0.667, full-continuation mass changes by $-0.260$ and energy by $+4.090$. Evaluation-only changes are $+5.786$ and $+1.202$. Lower total expenditure in the frozen controls reflects early exhaustion, not efficient persistence. This comparison establishes a benefit of reopening the combined learning-and-exploration procedure in some cases, without isolating parameter updates from corrective actions or their charges.

\begin{table}[htbp]
\centering\small
\caption{Content use and finite recovery from matched actual body states. Each row averages all nine cases (three shapes sharing each of three seed tapes), including exact-zero endpoints. Capture spans sixteen evaluation cycles. Changes and total paid cost span the full sixteen-cycle content continuation or thirty-two-cycle recovery continuation. The latter evaluates its final sixteen cycles; the first sixteen reopen learning only in the relearn branch. Cost includes movement, maintenance and memory. Rows from different protocols do not share the same starting states.}
\label{tab:ecology-continuations}
\begin{tabular}{lrrrrrr}
\toprule
Branch & Capture & $\Delta M$ & $\Delta E$ & Final $E$ & Paid cost & Zero/9\\
\midrule
Intact cue & 1.000 & 8.501 & $-4.386$ & 50.828 & 77.262 & 0\\
Reset cue & 0.542 & 2.700 & $-32.832$ & 22.382 & 71.669 & 0\\
\midrule
Remain frozen & 0.000 & $-34.118$ & $-9.539$ & 0.000 & 20.168 & 9\\
Relearn, then freeze & 0.667 & $-0.260$ & 4.090 & 13.629 & 96.621 & 3\\
\bottomrule
\end{tabular}
\end{table}


The same ecology also lets us separate a fixed temporal representation from a learned prediction. A second controller encodes permitted cue history with damped modes whose dynamics never learn. Only its action-conditioned linear readout is fitted, by recursive least squares after completed resource uptake. Across two prescribed shapes and two seeds, its mean capture in the sixteen frozen test cycles is 0.997 under stable predictive contingencies, but only 0.047 after reversal; mean test mass and energy changes in the latter condition are $-29.903$ and $-17.576$. Freezing the readout at reversal instead gives capture 0.0001, while a current-observation readout gives 0.349. A fixed last-pair heuristic reaches 0.997 in that reversal task by exploiting the stipulated complementary cue mapping; it is not a learned model of that mapping. These pilot comparisons distinguish information retained by a controller from information it can successfully revise. They neither establish the necessity of oscillatory modes over simpler cue retention nor show that changing tissue has learned a temporal representation.

An exploratory follow-up uses a completed prediction shortfall to request one probe selected by inverse-design leverage without resetting the learned weights. On the same four construction--seed cells and environmental tapes, mean reversal capture increases to 0.368, but mean test mass and energy still fall by 14.346 and 19.170. The follow-up is a response to the observed failure, not independent confirmation, and its trigger is prescribed rather than learned. All policies, physical losses, unpaid access and extinct paths remain in the source summaries. Both readout experiments reverse during training and freeze at cycle 32. A read-only audit records three, three, four and three probe requests in cycles 24--31 across the four reversal paths; two still have pending requests at the freeze boundary. The protocol disables both learning and probing thereafter. These are requests rather than guaranteed travel or new information, and their cutoff cannot be interpreted as a test of continuous online correction. Their scientific implication is a testable design requirement: preserving an informative event trace is insufficient unless action continues to obtain evidence that can correct a failed prediction.
